\section{Null space, principal directions, and SVD}

A concrete null vector is the cross product of the two rows of $\Ae$:
\begin{equation}
 \vect n=\begin{bmatrix}
 -\omegae^2\Lq\Lexc\\
 -\Rs\omegae\Lexc\\
 \Rs^2+\omegae^2\Ld\Lq
 \end{bmatrix},\qquad \Ae\vect n=\vect0.\label{eq:general-null}
\end{equation}
The slice-center line in \cref{eq:slice-center}, augmented with its excitation
coordinate, is exactly a scalar multiple of this vector.  For $R_s=0$ it
simplifies to
\begin{equation}
 \vect n\sim\begin{bmatrix}-\Lexc/\Ld\\0\\1\end{bmatrix}.
 \label{eq:ideal-null}
\end{equation}
For any state and scalar $t$,
\begin{equation}
 \vect u(\state+t\vect n)=\Ae\state+t\Ae\vect n=\Ae\state.
 \label{eq:null-invariance}
\end{equation}
Movement along the cylinder axis costs exactly no additional stator voltage in
this model.  This is the physical content of the zero eigenvalue.

\begin{figure}[tb]
\centering
\begin{tikzpicture}
\begin{axis}[publication 3d,width=.84\linewidth,height=7.4cm,
 xlabel={$i_d$},ylabel={$i_q$},zlabel={$i_e$},
 xmin=-2.7,xmax=2.7,ymin=-1.3,ymax=1.3,zmin=-2.1,zmax=2.1]
 \addplot3[surf,shader=interp,opacity=.34,draw=MidnightBlue!55,
  domain=0:360,y domain=-1.8:1.8,samples=37,samples y=9]
  ({-.72*y+1.05*cos(x)},{.72*sin(x)},{y});
 \addplot3[null direction,domain=-1.8:1.8,samples=2]
  ({-.72*x},{0},{x});
\end{axis}
\end{tikzpicture}
\caption{The rank-two stator-voltage level is an oblique elliptic cylinder.
The orange line is its null-space axis.}
\label{fig:voltage-cylinder}
\end{figure}

\begin{figure}[tb]
\centering
\begin{tikzpicture}
\begin{axis}[publication 3d,width=.75\linewidth,height=6.5cm,
 xlabel={$i_d$},ylabel={$i_q$},zlabel={$i_e$},
 xmin=-1.4,xmax=1.7,ymin=-1.3,ymax=1.4,zmin=-1.2,zmax=1.6]
 \addplot3[eigenvector] coordinates {(0,0,0) (1.25,.15,.25)};
 \addplot3[eigenvector] coordinates {(0,0,0) (-.1,1.05,.2)};
 \addplot3[null direction] coordinates {(0,0,0) (-.65,0,1.2)};
 \node[annotation] at (axis cs:1.2,.18,.35){$\vect v_1,\ \sigma_1>0$};
 \node[annotation] at (axis cs:-.1,1.05,.3){$\vect v_2,\ \sigma_2>0$};
 \node[annotation] at (axis cs:-.75,0,1.3){$\vect v_3,\ \sigma_3=0$};
\end{axis}
\end{tikzpicture}
\caption{Two principal directions have finite voltage gain; the third is the
null direction and gives an infinite cylinder axis.}
\label{fig:eesm-principal-axes}
\end{figure}


The full \gls{svd} can be written
\begin{equation}
 \Ae=\matr U
 \begin{bmatrix}\sigma_1&0&0\\0&\sigma_2&0\end{bmatrix}
 \matr V\trans,
 \qquad \sigma_1,\sigma_2>0.
\end{equation}
Then
\begin{equation}
 \Qu=\matr V\diag(\sigma_1^2,\sigma_2^2,0)\matr V\trans.
\end{equation}
The three right singular vectors are current-space principal directions.  The
first two produce voltage gains $\sigma_1$ and $\sigma_2$ and have cylinder
semi-axes
\begin{equation}
 a_1=\frac{\Umax}{\sigma_1},\qquad
 a_2=\frac{\Umax}{\sigma_2}.
\end{equation}
The third has $\sigma_3=0$ and therefore $a_3=\infty$: a compact ellipsoid is
impossible.  The rank argument, eigenvalue signature, null-space invariance,
and SVD all express the same geometry at different levels of detail.

\begin{figure}[tb]
  \centering
  \begin{tikzpicture}[
    node distance=9mm,
    box/.style={draw,rounded corners,align=center,minimum height=12mm,
      text width=.19\linewidth,fill=gray!6},
    flow/.style={-{Latex[length=2mm]},thick}
  ]
    \node[box] (input) {$\R^3$ current space\\three directions};
    \node[box,right=of input] (rotate) {$\matr V\trans$\\rotate to principal
      currents};
    \node[box,right=of rotate] (scale) {$[\matr\Sigma\ \vect0]$\\scale two,
      collapse one};
    \node[box,right=of scale] (output) {$\matr U$\\rotate in $\R^2$ voltage
      space};
    \draw[flow] (input)--(rotate);
    \draw[flow] (rotate)--(scale);
    \draw[flow] (scale)--(output);
    \node[below=5mm of scale,align=center,text width=.46\linewidth] {
      $\sigma_1,\sigma_2>0$: finite cross-section;\qquad
      $\sigma_3=0$: voltage-invisible cylinder axis};
  \end{tikzpicture}
  \caption{The EESM SVD as a sequence of operations. Rotations preserve shape
  and dimension; the rectangular scaling stage is where one input direction
  is removed. Traversing the map backward turns the voltage disk into an
  elliptic cylinder in current space.}
  \label{fig:eesm-svd-pipeline}
\end{figure}


\subsection{Five descriptions, one geometric fact}

\begin{center}
\small
\begin{tabularx}{\linewidth}{@{}p{.20\linewidth}X@{}}
\toprule
Description & Most useful interpretation\\
\midrule
Dimension and rank &
The fastest mathematical classification: three inputs, two independent
outputs, and therefore at least one invisible input direction.\\
Null space &
The clearest controller and physical interpretation: it gives the precise
current tradeoff that leaves stator voltage unchanged.\\
Eigenvalue signature &
The standard quadric language: two positive curvatures and one zero curvature
mean an elliptic cylinder rather than an ellipsoid.\\
SVD &
The best numerical algorithmic view: right singular vectors give current
directions, singular values give voltage gains and conditioning, and the zero
singular value detects the free axis robustly.\\
Stacked slices &
The most direct visualization: congruent fixed-excitation ellipses translate
along one line and assemble the cylinder.\\
\bottomrule
\end{tabularx}
\end{center}

These descriptions agree because
$\Qu=\Ae\trans\Ae$.  Rank counts visible directions, the null space identifies
the missing one, eigenvalues square the directional gains, the SVD retains
both input and output directions, and the slice picture reconstructs the same
set one excitation plane at a time.  None of the five is redundant: each
reveals information that another representation suppresses.

\paragraph{What if excitation were bounded?}
The voltage surface itself would still be a cylinder, but intersecting it with
$i_{e,\min}\leq i_e\leq i_{e,\max}$ would cut out a finite slab.  Boundedness
would come from the added excitation constraint, not from a change in the
rank-two voltage geometry.
